\subsection{Proposed Solution}

\subsubsection{Solution Architecture}
The \textbf{Iterator Pattern}~\cite{gof1994} decouples traversal algorithms from collection data structures by extracting navigation state into an external, dedicated object known as an \textbf{Iterator}. 

The pattern establishes two primary interface hierarchies~\cite{gof1994, refactoring_guru_iterator}:
\begin{enumerate}
    \item \textbf{Aggregate Interface (\texttt{Iterable}):} Defines a factory method (e.g., \texttt{createIterator()}) that supplies a compatible traversal cursor without revealing the internal container layout.
    \item \textbf{Iterator Interface:} Exposes a standard traversal contract—traditionally \texttt{hasNext()} and \texttt{next()}, or \texttt{first()}, \texttt{next()}, \texttt{isDone()}, and \texttt{currentItem()}.
\end{enumerate}

\begin{figure}[H]
    \centering
    \begin{tikzpicture}[
        node distance=2.2cm and 2.5cm,
        classbox/.style={
            draw, 
            rectangle, 
            minimum width=4.4cm, 
            minimum height=1.6cm, 
            align=center, 
            fill=blue!5, 
            font=\small\sffamily
        },
        arrow/.style={-latex, thick},
        dashedArrow/.style={-latex, thick, dashed},
        inherit/.style={-{Triangle[open, length=3mm, width=3mm]}, thick}
    ]

    % Aggregate Interface
    \node (aggInt) [classbox] {
        \textbf{\textit{<<Interface>>}}\\[1pt]
        \textbf{\textit{TrackAggregate}}\\[2pt]
        \rule{4.0cm}{0.4pt}\\[3pt]
        \footnotesize+ createIterator(): unique\_ptr<TrackIterator>
    };

    % Iterator Interface
    \node (iterInt) [classbox, right=2.0cm of aggInt] {
        \textbf{\textit{<<Interface>>}}\\[1pt]
        \textbf{\textit{TrackIterator}}\\[2pt]
        \rule{4.0cm}{0.4pt}\\[3pt]
        \footnotesize+ hasNext(): bool\\[1pt]
        \footnotesize+ next(): const Track\&
    };

    % Concrete Aggregate
    \node (aggConc) [classbox, below=1.8cm of aggInt, minimum height=2.2cm] {
        \textbf{PlaylistQueue}\\[2pt]
        \rule{4.0cm}{0.4pt}\\[3pt]
        \footnotesize-- tracks: vector<Track>\\[1pt]
        \rule{4.0cm}{0.4pt}\\[3pt]
        \footnotesize+ addTrack(t: Track): void\\[1pt]
        \footnotesize+ createIterator(): unique\_ptr<TrackIterator>
    };

    % Concrete Iterator
    \node (iterConc) [classbox, below=1.8cm of iterInt, minimum height=2.2cm] {
        \textbf{VectorTrackIterator}\\[2pt]
        \rule{4.0cm}{0.4pt}\\[3pt]
        \footnotesize-- items: const vector<Track>\&\\[1pt]
        \footnotesize-- cursor: size\_t = 0\\[1pt]
        \rule{4.0cm}{0.4pt}\\[3pt]
        \footnotesize+ hasNext(): bool\\[1pt]
        \footnotesize+ next(): const Track\&
    };

    % Client
    \node (client) [classbox, above=1.5cm of aggInt, xshift=3.2cm, minimum height=1.2cm, fill=green!5] {
        \textbf{Client / Consumer (\texttt{printPlaylist})}\\[2pt]
        \rule{6.5cm}{0.4pt}\\[3pt]
        \footnotesize+ printPlaylist(playlist: const TrackAggregate\&): void
    };

    % Connections
    \draw[inherit] (aggConc) -- (aggInt);
    \draw[inherit] (iterConc) -- (iterInt);
    \draw[dashedArrow] (aggConc) -- node[above, font=\scriptsize] {creates} (iterConc);
    \draw[dashedArrow] (iterConc) -- ++(-1.6, 0) -- node[below, font=\scriptsize, pos=0.5] {references container} ([yshift=-0.3cm]aggConc.east);
    \draw[arrow] (client.south -| aggInt.north) -- (aggInt.north);
    \draw[arrow] (client.south -| iterInt.north) -- (iterInt.north);

    \end{tikzpicture}
    \caption{UML Class Diagram for Concrete Media Playlist Iterator Implementation (Adapted from~\cite{viblo_iterator})}
    \label{fig:playlist_iterator_class_diagram}
\end{figure}

\subsubsection{C++ Implementation Example}
Below is the robust C++ implementation demonstrating the aggregate queue decoupled from its consumer via an iterator interface. Notice the defensive boundary checking inside \texttt{next()} to eliminate undefined behavior:

\begin{lstlisting}[language=C++, caption={C++ Implementation of Playlist Iterator Pattern}]
#include <iostream>
#include <memory>
#include <vector>
#include <string>
#include <stdexcept>

// Target data element
struct Track {
    std::string title;
    std::string artist;
};

// 1. Iterator Interface
class TrackIterator {
public:
    virtual ~TrackIterator() = default;
    virtual bool hasNext() const = 0;
    virtual const Track& next() = 0;
};

// 2. Aggregate Interface
class TrackAggregate {
public:
    virtual ~TrackAggregate() = default;
    virtual std::unique_ptr<TrackIterator> createIterator() const = 0;
};

// 3. Concrete Aggregate
class PlaylistQueue : public TrackAggregate {
private:
    std::vector<Track> tracks;

    // Concrete Iterator as an inner/friend class
    class VectorTrackIterator : public TrackIterator {
    private:
        const std::vector<Track>& items;
        size_t cursor = 0;

    public:
        explicit VectorTrackIterator(const std::vector<Track>& target)
            : items(target) {}

        bool hasNext() const override {
            return cursor < items.size();
        }

        const Track& next() override {
            if (!hasNext()) {
                throw std::out_of_range("VectorTrackIterator::next(): Iterator out of bounds.");
            }
            return items[cursor++];
        }
    };

public:
    void addTrack(const Track& track) {
        tracks.push_back(track);
    }

    std::unique_ptr<TrackIterator> createIterator() const override {
        return std::make_unique<VectorTrackIterator>(tracks);
    }
};

// 4. Client Usage (Uniform Traversal)
void printPlaylist(const TrackAggregate& playlist) {
    std::unique_ptr<TrackIterator> it = playlist.createIterator();
    int index = 1;
    while (it->hasNext()) {
        const Track& t = it->next();
        std::cout << index++ << ". " << t.title 
                  << " by " << t.artist << "\n";
    }
}
\end{lstlisting}